% =====================================================================================================
%  Radiology: Artificial Intelligence -- Original Research (focused clinical paper)
%  Every number in this file is read from one of:
%    docs/COUNTING_EVIDENCE.md         (all counting and naming arms, generated from the result JSONs)
%    docs/WRONG_LEVEL_COST_EVIDENCE.md (rates, litigation, fiducial cost; each claim carries its vote)
%    docs/WRONG_SITE_BURDEN_EVIDENCE.md(annual burden; extrapolations labelled)
%    docs/SWEEP_TIME_EVIDENCE.md       (C-arm move times, exposure times, competitor latency)
%    docs/SWEEP_PROTOCOL.md            (the protocol handed to the radiographer)
%  Nothing is typed from memory. Arms still in flight are marked [PENDING] with the arm's name.
% =====================================================================================================
\documentclass[11pt,letterpaper]{article}
\usepackage[margin=1in]{geometry}
\usepackage[T1]{fontenc}
\usepackage{lmodern}
\usepackage{helvet}
\renewcommand{\familydefault}{\sfdefault}
\usepackage{setspace}\setstretch{1.15}
\usepackage{graphicx,booktabs,array,amsmath,textcomp,microtype}
\usepackage[labelfont=bf,font=small]{caption}
\usepackage[numbers,sort&compress]{natbib}
\usepackage[hidelinks]{hyperref}

\begin{document}

% ------------------------------------------------------------------------------------------- TITLE
\begin{center}
{\Large\bfseries Counting Made Exact: Registration-Free Confirmation of the Operative\\[2pt]
Vertebral Level from a Short Fluoroscopic Sweep\par}
\vspace{0.9em}
{\normalsize Gregory Schwing, MD, PhD\par}
{\small Department of Neurosurgery, Wayne State University School of Medicine /
Detroit Medical Center, Detroit, MI\par}
{\small \texttt{go2432@wayne.edu}\par}
\end{center}

\vspace{0.5em}
\noindent\textbf{Article type:} Original Research. \textbf{Main text:} about 3{,}200 words.
\textbf{Tables:} 3. \textbf{Figures:} 0. \textbf{Funding:} none declared.
\textbf{Conflicts of interest:} the author has a commercial interest in a level-confirmation product
under development.

\medskip
\noindent\textbf{Summary statement.} Linking a short series of overlapping fluoroscopic exposures into
one chain, with no learned component in the count, placed the mark on the correct vertebra in 834 of 834
walks from C2 on 76 whole-spine studies and in 99.66\,\% of 6{,}255 walks from the sacrum on 801
abdominopelvic studies, using the annotation's vertebra boxes.

% ------------------------------------------------------------------------------------------ ABSTRACT
\section*{Abstract}

\noindent\textbf{Background.} Wrong-level spine surgery persists under the Universal Protocol because
the protocol verifies a site that has been \emph{marked}, and a vertebral level has no mark. The level is
counted in theatre across several limited fluoroscopic fields, and each join between fields can be one
body off.

\noindent\textbf{Purpose.} To measure whether counting across overlapping exposures can be made exact
without any learned component, and to locate the residual risk.

\noindent\textbf{Materials and Methods.} Retrospective study of public data. Sweeps were simulated from
labelled volumes and radiographs by the protocol a radiographer is handed: put the anchor (sacrum or C2)
at the edge of the field, expose, move one field minus two vertebral bodies, expose again. Consecutive
frames were linked by normalised cross-correlation restricted to the step the protocol commanded, which is
what prevents the periodic column from aliasing one body. The target is the $k$-th node from the anchor;
no vertebra is named and no network is loaded for the count. A walk is \emph{exact} when the mark lands on
the body truly $k$ from the anchor, scored among walks that reached the target before the image ran out;
both denominators are reported. \emph{All headline figures use the annotation's vertebra boxes}; arms
using a detector's own boxes are labelled as such.

\noindent\textbf{Results.} From C2 at three bodies of overlap, 834 of 834 walks on 76 whole-spine studies
were exact at 2.24 exposures; at two bodies, 832 of 834. On 204 TotalSegmentator studies, 99.70\,\%
(2{,}646/2{,}654 reached). From the sacrum on all 801 abdominopelvic studies, 99.66\,\% (2-body overlap,
6{,}255 reached) and 99.76\,\% (3-body, 6{,}229 reached), at about two exposures. On real radiographs with
the deployed film detector's own boxes, 95 of 95 walks on 16 spines were exact, and from C2 with the
annotation's boxes, 110 of 110 on 24 spines. The residual risk is in the boxes: on the same 79 CT studies,
true boxes gave 100\,\% and a render-trained detector 30.06\,\%, with the node the count started from being
the true sacrum on 42.44\,\% of walks.

\noindent\textbf{Conclusion.} Counting across overlapping fields is exact and needs no network; the
uncertainty lives in finding the anchor, and in converting an ordinal into a name.

% -------------------------------------------------------------------------------------- INTRODUCTION
\section*{Introduction}

The patient is prone and draped, the fascia is open, and the surgeon has placed a spinal needle in what
should be the target interspace. The C-arm swings in for a lateral shot. On the monitor are three or four
vertebral bodies and nothing else: no sacrum, no ribs, no odontoid. The surgeon counts back to the last
exposure, which held three or four other bodies, and to the one before that, and arrives at a name for the
body the needle is touching. Nothing in the picture confirms that name. If any join between exposures was
made one body off, the needle is in the wrong interspace and every subsequent step of the operation is
performed there.

This is not a rare manoeuvre performed badly. It is the standard technique, unchanged for decades: needles
placed every few levels and counted cephalad from the sacrum on lateral fluoroscopy, with surgeons
disagreeing even on which structure to count---the pedicle for 67\,\% of respondents to a national survey
and the spinous process for 49\,\% \citep{strong2021,mayer2008}. Half of responding neurosurgeons in two
independent surveys report at least one wrong-level operation in their career \citep{mody2008,groff2013}.
Self-reported incidence is one in 3{,}110 procedures \citep{mody2008}; a prospectively collected registry
of 34{,}544 degenerative spine operations recorded 0.12\,\% (95\,\% CI 0.09--0.16), with the thoracic spine
highest at 0.42\,\% (CI 0.16--1.08), three times the lumbar and eight times the cervical rate, and
intraoperative disorientation the cause in 81\,\% of events \citep{furuya2026}.

The Universal Protocol does not close this. Its three elements---verification, site marking, time-out---all
assume the site can be identified on the patient before incision. For a limb that is true. For a vertebra
it is not: there is no skin mark that names L4. Chart review of 13 insured wrong-site cases judged that the
protocol, applied under ideal conditions, would have prevented eight of them, 62\,\% \citep{kwaan2006}; the
single study to assess site-verification protocols in the spine found they failed to prevent about
one-third of cases \citep{longo2012}; and the literature's recommended safeguard, imaging after exposure
against a marked fixed structure \citep{devine2010}, is itself another count. Wrong-site procedures
continued at a steady rate through the protocol era, with significant harm in 35.5\,\% of them
\citep{stahel2010}. The Joint Commission received 127 reports of wrong surgery in 2024, nine of them
wrong-level and seven of those spinal, from a voluntary system it states captures a small fraction of
events \citep{jc2024}; NHS England recorded 185 wrong-site never events in a single year \citep{nhse2025}.

What it costs to avoid counting is measurable, and it is paid in the thoracic spine. One quaternary centre
placed 179 preoperative percutaneous gold fiducial markers over 38 months, 176 of them thoracic, at a mean
of 51.9 minutes and 183\,mGy each, purely to mark a level for the surgeon; all 179 operations were at the
correct level \citep{siminski2023}. At a hospital operating-room cost of \$33.78 per minute (median,
FY2014, before physician fees) that is roughly \$1{,}750--\$1{,}950 of staffed procedural time per patient
and about 155 hours across the series \citep{childers2018}---a valuation by proxy, since the marker is
placed in a fluoroscopy or CT suite rather than a theatre. A second series puts the CT intervention itself
at 14.3 minutes from scout to control scan \citep{keil2024}. With the marker in place, intraoperative
localisation falls from about 15 minutes to about three \citep{mongardi2021,upadhyaya2012}. The profession
has decided, in the region where counting is longest, that a second invasive procedure scheduled in advance
is worth it.

The obvious machine response is to name the vertebra in the picture. It is the wrong response, for a reason
visible in the operating room rather than in a validation set. The surgeon is not attempting a
classification and does not want one. The plan says ``L4--L5'', and the surgeon reaches it by counting from
a landmark; a machine that asserts ``this is L4'' from morphology offers a second opinion of the same kind
the surgeon already distrusts, formed from a picture in which adjacent vertebrae genuinely resemble one
another. Worse, the name is not a property of the bone. A lumbosacral or thoracolumbar transitional
vertebra changes how many bodies the column has, so two spines with identical local anatomy carry different
names---and transitional spines are the population in which wrong-level surgery occurs. Registration-based
systems side-step naming by importing it: they solve a pose between a preoperative CT and the radiograph
and transfer the name a human read off the volume \citep{levelcheck2012,otake2015}. That requires a CT the
patient may not have, and a miscounted CT propagates onto the radiograph faithfully.

The intervention the operation actually needs is therefore not a classifier. It is a mechanism that makes
the surgeon's own count exact across fields. This study measures one: overlapping exposures from an anchor,
linked into a single chain, with the target defined as the $k$-th body from the anchor. No level name is
formed and no network is loaded for the count.

% ------------------------------------------------------------------------------------------- METHODS
\section*{Materials and Methods}

This retrospective study used public, de-identified data only; institutional review determined it not to be
human-subjects research.

\subsection*{The sweep, as a radiographer performs it}
The protocol is written for the person at the C-arm \citep{strong2021}. Find the anchor and put it at the
edge of the field: the sacral promontory at the bottom for lumbar and thoracolumbar work, C2 at the top for
cervical and upper thoracic work. Expose. Move the C-arm toward the level of interest by one field minus
two vertebral bodies, so that two bodies of the previous picture remain in view. Expose. Repeat until the
level of interest is in the field, and stop---the sweep runs from the anchor to the work, not along the
whole spine. The field size does not change during a sweep. The last move is a full move or none: a final
frame that advances less than half a body is nearly the same picture twice, and cannot be distinguished
from a step one whole body longer.

\subsection*{Linking, and why the restriction matters}
Consecutive frames are linked by normalised cross-correlation computed on the overlap band alone, never by
a global transform, and the search is restricted to within 0.6 of a vertebral body of the step the protocol
commanded. The restriction is the mechanism, not a tuning choice. The column is nearly periodic; an
unrestricted correlation between two lumbar frames has a second peak one body away that is almost as strong
as the correct one, and locking onto it is precisely the one-body-off error the surgeon makes. Because the
protocol states the step, the software is asked only to refine a displacement it was told, which is what
removes the aliasing.

\subsection*{The count}
The linked frames form one chain of vertebrae. Each node is identified by its ordinal from the anchor, and
the target is the $k$-th node. Nothing is named, no alphabet is chosen, and no learned component is
loaded. Two consequences follow. Variant anatomy cannot enter, because an ordinal is well defined on a
column with five or six lumbar bodies alike. And because nothing is fitted, no spine has to be withheld:
every study carrying the anchor is used, and there is no training-distribution question for the counting
step.

\subsection*{Data}
Sweeps were cut from labelled long images, because serial sweeps of one patient exist in no public
collection. Counting from the sacrum used all 801 abdominopelvic CT studies of CTSpinoPelvic1K
\citep{ctsp1k} and the 522 TotalSegmentator studies that reach the sacrum; counting from C2 used 76
whole-spine VerSe studies \citep{verse2021} and the 204 TotalSegmentator studies that hold C2.
TotalSegmentator is routine clinical CT from another institution on which no model in this work has been
trained. Real radiographs contributed 16 spines that reach the sacrum and 24 in which C2 is visible; these
counts are small and are stated beside every percentage that rests on them.

\subsection*{Scoring, and the two denominators}
A walk is \emph{exact} when the mark lands on the body that is truly $k$ from the anchor. Walks are scored
among those that \emph{reached} the target---that is, in which the target body entered a field before the
image ran out. A walk that does not reach is a limitation of the field, not a counting error, and both
denominators are reported in Table~\ref{tab:count}. \emph{Anchor held} is the share of walks in which the
node the count started from really was the anchor; it is a diagnostic and never a gate.

\subsection*{Boxes}
Every headline figure uses the vertebra boxes supplied by the annotation, so that the only thing that can
fail is the link. Two further conditions were run to locate the residual risk. On 79 CT studies held out
from a render-trained detector, the same sweeps were run with that detector's boxes, taking the anchor to
be the terminal node of the chain. On the 16 real radiographs, the sweeps were run with the deployed film
detector's own boxes. Arms that identify the anchor by \emph{class} rather than by position are still in
flight and are marked [PENDING] in Table~\ref{tab:boxes}: \textsc{sacrum, CTSP, detector boxes,
anchor=class}; \textsc{sacrum o3, CTSP, detector boxes, anchor=class}; \textsc{sacrum, films, film
detector, anchor=class}; and a C2-class anchor detector on real film.

\subsection*{Naming, for contrast}
A separate arm names the chain instead of counting it: a ResNet-18 encoder over fixed-scale vertebra crops
with a recurrent head over the chain, decoded under the constraint that the chain is consecutive
\citep{he2016,schwing2026ctx,schwing2026mixed}. It is reported here only to contrast where uncertainty
lives, and its per-frame and per-walk compute is in Table~\ref{tab:time}.

\subsection*{Not measured here}
No hypothesis test was applied. The sweeps are simulated from single long images: each frame is degraded
independently, so no two frames share pixels exactly, but parallax between exposures taken at different
C-arm positions and patient motion during a sweep are not reproduced. Both act on the linking and neither
on the count. A needle-confirmation arm---re-exposing with a metallic instrument in the field and placing
the instrument tip on the chain---is specified but not measured [PENDING: needle confirmation].

% ------------------------------------------------------------------------------------------- RESULTS
\section*{Results}

\subsection*{Counting is exact, at two to three exposures}
Table~\ref{tab:count} gives every arm. Counting down from C2 on 76 whole-spine studies was exact on 834 of
834 walks at three bodies of overlap, from a mean of 2.24 exposures, and on 832 of 834 at two bodies, from
1.99; both failures fell in one spine. On 204 TotalSegmentator studies that hold C2, 2{,}646 of the 2{,}654
walks that reached their target were exact (99.70\,\%) at two bodies of overlap and 99.62\,\% at three, the
misses confined to two spines in each case.

Counting up from the sacrum on all 801 abdominopelvic studies was exact on 99.66\,\% of the 6{,}255 reached
walks at two bodies of overlap, from 1.98 exposures, and 99.76\,\% of 6{,}229 at three, from 2.42; the
errors sat in six and five spines respectively. On the 522 TotalSegmentator studies that reach the sacrum
the figure was 99.17\,\% of 7{,}228 reached walks from 2.96 exposures, with 60 errored walks in six spines.
What does not reach is a limitation of the field rather than of the count: 97 of 6{,}352 walks on the
abdominopelvic studies and 168 of 7{,}396 on TotalSegmentator ended before the target appeared.

The residual errors are regional and they sit where the anchor is far away. From the sacrum on the
abdominopelvic studies, lumbar targets were exact 99.73\,\% of the time and thoracic targets 95.45\,\%; from
the sacrum on TotalSegmentator, lumbar 99.33\,\%, thoracic 97.41\,\% and cervical 83.84\,\%. From C2 the
gradient runs the other way: cervical 100\,\% and thoracic 99.46\,\% on VerSe, lumbar 97.87\,\% and sacral
97.67\,\% on TotalSegmentator. A sweep from the near anchor is the operating instruction that follows.

On real radiographs, all 95 walks on 16 spines reaching the sacrum were exact, from 1.61 exposures, and all
110 walks on 24 spines holding C2, from 1.15 exposures. Sixteen and 24 spines are thin, and the number is
stated beside the percentage for that reason.

\subsection*{The residual risk is in the boxes, and specifically in the anchor}
Table~\ref{tab:boxes} isolates it on one set of cases. On 79 CT studies held out from a render-trained
detector, the annotation's boxes gave 618 of 618 reached walks exact. The same sweeps run with that
detector's boxes gave 30.06\,\%, with 435 errored walks spread across all 79 spines; at three bodies of
overlap, 28.66\,\%. More overlap does not rescue bad boxes.

The diagnostic says why. Taking the anchor to be the terminal node of the chain, that node really was the
sacrum on 42.44\,\% of walks at two bodies of overlap and 39.77\,\% at three. A missed or spurious box at
the bottom of the chain shifts every ordinal above it by one, which is the surgeon's own error
reintroduced by a machine, delivered with the confidence of a computed result. Lumbar targets fell to
39.29\,\% and thoracic to 13.60\,\% under these boxes---the further from the anchor, the longer the shift
persists.

The failure is not universal to detectors. With the deployed film detector's own boxes on the 16 real
radiographs, all 95 walks were exact from 1.47 exposures, with the anchor diagnostic at 93.68\,\%. The
render-trained detector is the arm that fails, and it fails in the render domain. The remedy the arms
in flight test is to identify the anchor by class and refuse the walk when no anchor is found, rather than
assume the end of the chain is one [PENDING: anchor-by-class arms].

\subsection*{Naming is where the uncertainty is}
Run as a naming problem instead, the same chains were named correctly on 96.25\,\% of 950 walks on the
abdominopelvic studies (lumbar 94.98\,\%, thoracic 93.75\,\%) and 94.69\,\% of 7{,}405 walks on
TotalSegmentator (cervical 80.35\,\%, lumbar 96.21\,\%, thoracic 91.58\,\%). Those are good figures for a
classifier and they are three to five points worse than the count on the same corpora, which is the result
this paper is built around. The walk sets are not identical---950 and 7{,}405 walks against 6{,}352 and
7{,}396---so the comparison is between arms on the same studies, not between two readings of one walk.

\subsection*{Time}
Table~\ref{tab:time} places the compute against the manoeuvre. The encoder runs at 17--18\,ms per frame,
linking at 0.29\,s per walk on the abdominopelvic sweeps and 0.10\,s on TotalSegmentator, and decoding at
0.2\,ms. A C-arm move costs 36--49\,s in the only published time-motion measurements
\citep{matthews2007,desilva2018}, and physical manipulation rather than beam-on time dominates that
figure. At one to three exposures per walk, the compute is not the cost of this method and is not the
argument for it.

% ---------------------------------------------------------------------------------------- DISCUSSION
\section*{Discussion}

\subsection*{The division of labour}
Counting is exact and needs no network. Naming is where uncertainty lives, and it lives there for a reason
that no additional data removes: converting an ordinal into a name requires an alphabet, and transitional
anatomy changes the alphabet. A spine with six lumbar bodies and a spine with five can present identical
local morphology at the operative level and carry different names for the same bone. The network's proper
job is therefore not to identify each vertebra from its shape. It is to decide which alphabet this patient
has---and when the operative plan is stated relative to the anchor, that decision does not have to be made
at all.

This reframes what the machine is for. It does not tell the surgeon what the level is. It performs the
surgeon's own count, across the same fields, without the join errors, and it declines when it cannot see
the anchor. The output is auditable by eye, because the chain of overlapping images is on the screen and a
reader can check a link by looking at the bodies two frames share. That is a different object from a
classifier's posterior, which cannot be checked by inspection at all.

\subsection*{Where it fits in the Universal Protocol}
The protocol's gap in spine surgery is structural: it verifies a mark and the spine has none. The sweep
supplies one. Verification enters the target level as an ordinal from a named anchor. Site marking becomes
the sweep itself, which places a digital fiducial on the target body in the live image. The time-out
compares that fiducial's count against the surgeon's own, and a disagreement is a stop rather than an
override. This is an addition to the protocol, not a replacement for anyone's judgement, and its primary
quality metric is the disagreement rate---a number no published study of intraoperative localisation
reports, because the localisation \emph{process} has never been measured as a process.

\subsection*{Confirming with an instrument}
The same mechanism extends to the confirmation a surgeon already improvises. With the needle or a
radiopaque instrument in the field, one further exposure from the same or a slightly different angle can be
matched onto the chain, and the body the instrument tip lies over carries its ordinal from the anchor. This
reproduces what optical navigation provides---an instrument localised on the patient's own anatomy---with a
C-arm the hospital already owns and an embedded computer, and with no tracker, no preoperative CT and no
line of sight to maintain. It has not been measured. The arm is specified and its results are not in this
paper [PENDING: needle confirmation].

\subsection*{Deployment}
The case for deployment is not speed. The compute is negligible in both designs (Table~\ref{tab:time})
against 36--49\,s per C-arm move, and registration-based labelling has been reported at 4.7\,s per image on
a 2012 GPU and 17--72\,s in an operating-room deployment \citep{levelcheck2012,manbachi2018}. The case is
that the count needs one learned component rather than two. Naming requires a detector to find the
vertebrae and an encoder to name them; counting requires only the detector, and the anchor test it must
pass. It needs no preoperative CT, which is what confines registration systems to the patients who have
one, and no optical tracking hardware. A comparison against the capital cost of navigation suites is not
made here: no verified figure for that cost was among the sources used for this paper, and one is not
invented.

\subsection*{What would be saved is bounded, not known}
The burden this addresses cannot be stated as a worldwide or annual dollar figure from the evidence
available. The often-quoted \$1.3 billion is 21 years of cumulative malpractice payments for all surgical
never events, about \$62 million a year, of which wrong-site is roughly a quarter \citep{mehtsun2013}; the
mean closed wrong-site claim in a 2013--2020 carrier series was \$136{,}453, with spine and disc surgery
the commonest procedure type \citep{tan2023}; the median indemnity for a surgical never event reported to
the National Practitioner Data Bank is \$75{,}220 (IQR \$28{,}847--\$215{,}003) \citep{bettis2026}; 978 NHS
spinal-surgery claims over five years were valued at \pounds535.5 million \citep{machin2018}. Applying the
two verified per-procedure rates to a United States volume of 1.3--1.6 million spine operations a year
gives, as derived arithmetic on an unverified denominator, roughly 400--1{,}900 wrong-level operations and
\$8--37 million of indemnity a year. Seventeen percent of self-reported wrong-level events led to legal
action or settlement \citep{mody2008}, 55 of 68 litigated wrong-disc-level cases ended in a payment and
surgeons prevailed in 42\,\% of those a jury decided \citep{goodkin2004}, and nearly one in five Spine
Section respondents has been the subject of such a case \citep{groff2013}. No verified source gives a
per-event hospital cost, and no defensible worldwide figure exists.

\subsection*{Limitations}
Every headline counting figure uses the annotation's vertebra boxes, so it measures the linker and not the
whole system; the detector-box arms are what measure the whole system, and one of them fails badly. The
real-film arms rest on 16 and 24 spines. The sweeps are simulated: parallax between exposures at different
C-arm positions and patient motion within a sweep are not reproduced, and both make linking harder. The
anchor-by-class arms that would answer the central failure are in flight and unreported, as is the
needle-confirmation arm and the C2-class anchor detector for real film. The naming arms use models fitted
on some of the same collections and are reported for contrast only. No reader study has been run: the
honest comparison for a decision-support tool is against the count a surgeon would otherwise make, and that
comparison is not in this paper.

\subsection*{Conclusion}
The operation does not need a machine that knows what a vertebra is called. It needs the surgeon's count to
be exact across the fields the C-arm can show, and that is achievable with no learned component in the
count: 834 of 834 walks from C2 on 76 whole-spine studies, 99.66\,\% from the sacrum on 801 abdominopelvic
studies, at two to three exposures, with the annotation's boxes. What remains uncertain is finding the
anchor, which is a detector problem with a safe failure mode---refuse---and naming, which is an alphabet
problem that the plan can be written to avoid.

% -------------------------------------------------------------------------------------------- TABLES
\clearpage
\section*{Tables}

\begin{table}[htbp]\centering
\caption{Counting arms, with the annotation's vertebra boxes. \emph{Reached} is the number of walks in
which the target entered a field before the image ran out; \emph{exact} is scored on those.
\emph{Anchor held} is a diagnostic, never a gate.}
\label{tab:count}
\small
\resizebox{\textwidth}{!}{%
\begin{tabular}{llrrrrrr}
\toprule
Anchor, overlap & Corpus & Spines & Reached/walks & Exact & Exposures & Anchor & Wrong (spines) \\
\midrule
C2, 2 bodies     & VerSe             &  76 & 834/834     & 99.76\,\% & 1.99 & 100\,\%    &  2 (1) \\
C2, 3 bodies     & VerSe             &  76 & 834/834     & 100\,\%   & 2.24 & 100\,\%    &  0 (0) \\
C2, 2 bodies     & TotalSegmentator  & 204 & 2{,}654/2{,}656 & 99.70\,\% & 2.05 & 100\,\% &  8 (2) \\
C2, 3 bodies     & TotalSegmentator  & 204 & 2{,}650/2{,}656 & 99.62\,\% & 2.27 & 100\,\% & 10 (2) \\
Sacrum, 2 bodies & CTSpinoPelvic1K   & 801 & 6{,}255/6{,}352 & 99.66\,\% & 1.98 & 100\,\% & 21 (6) \\
Sacrum, 3 bodies & CTSpinoPelvic1K   & 801 & 6{,}229/6{,}352 & 99.76\,\% & 2.42 & 100\,\% & 15 (5) \\
Sacrum, 2 bodies & TotalSegmentator  & 522 & 7{,}228/7{,}396 & 99.17\,\% & 2.96 & 99.35\,\% & 60 (6) \\
Sacrum, 2 bodies & Real radiographs  &  16 & 95/95       & 100\,\%   & 1.61 & 100\,\%    &  0 (0) \\
C2, 2 bodies     & Real radiographs  &  24 & 110/110     & 100\,\%   & 1.15 & 100\,\%    &  0 (0) \\
\bottomrule
\end{tabular}}
\end{table}

\begin{table}[htbp]\centering
\caption{The same cases under three box conditions. The count is unchanged; only the source of the
vertebra boxes differs. The render-trained detector loses the anchor, and every ordinal above it shifts.}
\label{tab:boxes}
\small
\resizebox{\textwidth}{!}{%
\begin{tabular}{llrrrrr}
\toprule
Cases & Boxes, anchor rule & Spines & Reached/walks & Exact & Exposures & Anchor held \\
\midrule
CT, detector-held-out & Annotation                      & 79 & 618/625 & 100\,\%    & 1.98 & 100\,\% \\
CT, same cases        & Detector, anchor = terminal node & 79 & 622/625 & 30.06\,\% & 1.72 & 42.44\,\% \\
CT, same cases (3-body) & Detector, anchor = terminal node & 79 & 621/625 & 28.66\,\% & 1.91 & 39.77\,\% \\
CT, same cases        & Detector, anchor = class         & \multicolumn{5}{l}{[PENDING]} \\
Real film             & Film detector, anchor = terminal node & 16 & 95/95 & 100\,\% & 1.47 & 93.68\,\% \\
Real film             & Film detector, anchor = class    & \multicolumn{5}{l}{[PENDING]} \\
\bottomrule
\end{tabular}}
\end{table}

\begin{table}[htbp]\centering
\caption{Where the time goes. Measurements of this system are marked; the rest are published figures for
the manoeuvre and for the alternatives.}
\label{tab:time}
\small
\begin{tabular}{p{7.6cm}p{4.0cm}p{3.4cm}}
\toprule
Quantity & Value & Source \\
\midrule
C-arm repositioning, per move & 46\,s (CI 41--51) and 49\,s (CI 44--53); 36.4 $\pm$ 33.2\,s &
\citet{matthews2007}; \citet{desilva2018} \\
Beam-on, per spot exposure & about 0.5\,s (phantom) & \citet{desilva2018} \\
Exposures per walk & 1.15--2.96 & this study \\
Encoder, per frame & 17--18\,ms & this study \\
Linking, per walk & 0.29\,s; 0.10\,s & this study \\
Decoding, per walk & 0.2\,ms & this study \\
Registration-based labelling, per image & 4.7\,s (2012 GPU); 17--72\,s in theatre &
\citet{levelcheck2012}; \citet{manbachi2018} \\
Thoracic localisation by fluoroscopic counting & about 15\,min & \citet{mongardi2021,upadhyaya2012} \\
The same with a preoperative fiducial & about 3\,min, plus 14.3\,min of CT intervention and a 51.9\,min,
183\,mGy placement & \citet{mongardi2021}; \citet{keil2024}; \citet{siminski2023} \\
\bottomrule
\end{tabular}
\end{table}

\clearpage
\bibliographystyle{unsrtnat}
\bibliography{../../rsnai/refs}

\end{document}
